\documentclass[11pt]{article}
\usepackage[margin=0.8in]{geometry}
\usepackage{amsmath,microtype}
\newcommand{\guideheading}[1]{\par\medskip\noindent{\large\bfseries #1}\par\smallskip}
\begin{document}
\begin{center}
{\Large Guide: What the Affine Transfer Would Prove}\par\medskip
Collatz proof project\quad---\quad September 5, 2026
\end{center}
\guideheading{The result}
Lean now verifies that the full Collatz conjecture is equivalent to this
single transfer statement:
\[
\begin{gathered}
z\equiv2\pmod3\text{ and }z\text{ reaches one}\\
\Longrightarrow\quad9z+2\text{ reaches one}.
\end{gathered}
\]
Equivalently, reaching one must transfer from $3u+2$ to $27u+20$ for
every natural $u$. The implication itself has not been proved. This is
an exact reformulation of the goal, not a completed proof.

\guideheading{Why this pair appears}
Compare an odd seed $n$ with the smaller seed $n-1$. The previous
classical odd-run result gives two possibilities. In one case they
merge. In the other, their later values are $9z+2$ and $z$.
After excluding the short cases that already descend, the smaller
endpoint $z$ is always two modulo three.

If every smaller positive seed reaches one, then $n-1$ and its later
value $z$ do too. The proposed transfer would make $9z+2$ reach one,
and therefore make $n$ reach one. Lean checks this argument by strong
induction, including the even seeds and the short odd-run cases.

\guideheading{Boundedness is a different version}
Replacing ``reaches one'' with ``has a bounded orbit'' gives an equivalent
reformulation of nondivergence. In fact, if any unbounded orbit exists,
Lean proves that there must be a bounded $z\equiv2\pmod3$ whose partner
$9z+2$ is unbounded.

Boundedness alone allows nontrivial cycles. The reaching-one version
addresses them as well, so it matches the full original goal. We should
not treat the two transfer statements as interchangeable without proof.

\guideheading{What this changes in the research plan}
The affine transfer cannot be assumed as a routine stability property.
Even its restriction to one residue class contains the full unresolved
claim. The useful output is a precise target and a proof that it would
suffice. A future argument must explain why reaching one transfers
between these partners for every allowed input.

Finite examples cannot establish that statement: both members may
converge in every tested pair while the universal implication remains
unproved. The new Lean proofs use no experimental premise and do not
assert any progress on establishing the transfer itself.

\guideheading{Where to check the proofs}
\begingroup\raggedright
The main theorem is \texttt{collatz\_iff\_restricted\_affine\_transfer}
in \texttt{lean/Collatz/AffineConvergence.lean}. The boundedness
equivalences and failure-pair construction are in
\texttt{lean/Collatz/AffineBoundedness.lean}. Build these modules from
\texttt{lean/}. The repository theorem audit checks their statements
and axiom dependencies. The technical paper gives the proof structure
and relates it to the classical odd-run identities.
\par\endgroup
\end{document}
